\begin{table}[!htbp]
\centering
\caption{Returns relative to size screened comparison securities}\label{tab:comparators}
\small
\setlength{\tabcolsep}{5pt}
\begin{tabular}{llrrrlr}
\toprule
Type & Period & Events & Controls & Coefficient & 95\% interval & $p$ \\
\midrule
Addition & 2007--2016 & 78 & 787 & 3.90 & [0.10, 8.07] & 0.046 \\
Addition & 2017--2026 & 159 & 1609 & 1.53 & [0.36, 2.68] & 0.015 \\
Deletion & 2007--2016 & 68 & 415 & -3.24 & [-9.27, 1.17] & 0.210 \\
Deletion & 2017--2026 & 130 & 833 & -1.40 & [-2.75, 0.03] & 0.053 \\
\bottomrule
\end{tabular}
\begin{minipage}{\linewidth}\footnotesize\vspace{5pt}
The coefficient is on an event indicator in a regression of CAR $[AD-1,ED+1]$ on that indicator, pre-announcement returns, log market capitalization and review fixed effects. Comparison observations are security-review observations, not unique firms. The inherited size screens do not reconstruct the Nasdaq inclusion threshold and are not a regression discontinuity design. The pre-return control includes AD-1, which also enters the outcome; this overlap is an additional reason to treat the coefficients as descriptive adjustments. Intervals and $p$-values use 9,999 restricted wild cluster draws by review. Source: near\_additions/near\_additions\_results.csv.
\end{minipage}
\end{table}
